PostgreSQL is the selected central relational store; exact DDL, indexes and
transactional enforcement still require implementation and verification.
{\readabletable
\begin{longtable}{L{0.28}L{0.72}}
\caption{Proposed physical data representations.}\label{tab:physical-model}\\
\toprule\rowcolor{TableHeader}
\textbf{Family} & \textbf{Proposed representation} \\
\midrule\endfirsthead
\toprule\rowcolor{TableHeader}
\textbf{Family} & \textbf{Proposed representation} \\
\midrule\endhead
\multicolumn{2}{r}{\footnotesize\itshape Continued on next page}\\\endfoot
\bottomrule\endlastfoot
Identity/scope & UUID primary and operation ids; composite scope keys support same-workspace references. Integer revision counters support concurrency. \\
\midrule
Time & Full event/signing/acceptance instants use timestamptz; date for work/campaign/expiry dates. Workspace business\_timezone stores a validated timezone identifier. This type preserves the instant but not the original timezone/offset, so preserve capture-time IANA timezone metadata separately. Device timezone governs entry/display; the confirmed business timezone governs shared working-date assignment and filters. Minute-level input does not reduce timestamps to time-of-day values. \\
\midrule
Quantities/money & numeric for accepted assigned kg, other weights, areas, yield and amounts. Monetary amounts use two decimal places; weights/areas preserve source precision and calculations retain full stored values. The team verifies database precision and range constraints. \\
\midrule
Source/configuration & Schema-validated jsonb for other reported weights, invitation configuration, template fields and reviewed generation inputs. Core associations remain FKs. \\
\midrule
Files & Private object storage for source, template, signature and rendition bytes; StoredFile metadata and server-side format/size/access validation. Upload limit: 10~MB = 10,000,000 bytes; JPEG/JPG share the same accepted image format. \\
\end{longtable}}

The instant/timezone distinction follows the
\href{https://www.postgresql.org/docs/18/datatype-datetime.html\#DATATYPE-DATETIME-INPUT-TIMESTAMPS}{PostgreSQL date/time documentation}.

\paragraph{Essential keys and indexes.}
Use unique membership, coverage/assignment, attendance, plan, current-document-source,
worker/version signature and ticket-valuation keys. Unique delivery references
enforce one ticket/analysis; a unique version/worker key identifies a rendition.
Index campaign/area/date queries, worker intervals, document worker/version links and
attachment targets. Owner and whole-property/parcel uniqueness need explicit
partial/NULL handling; minimum-child counts and dependency checks are transactional.

Enforce planned and actual same-worker overlaps separately through exclusion or
equivalent serialised validation. Serialise execution reservations across a property
and its parcels, including start/reopening and coverage changes. Unlocked
read-then-insert checks or ordinary indexes alone are insufficient.

Store required official identification as text, preserving leading zeroes; use a
non-unique workspace/\allowbreak{}type/\allowbreak{}country/\allowbreak{}normalised-number lookup index. Worker and
operation UUIDs support client creation. Campaign lifecycle values are
\texttt{planned}, \texttt{in\_progress} and \texttt{closed}.

The physical design also requires an operation-outcome ledger, authorised
change/deletion feed and protected owner-review recovery store. Preserve template
snapshots and retained signed/cancelled-record references. Local prepared subsets,
outbox and file metadata need not mirror every central table. Local authorisation
metadata stores online-validation and expiry instants; its approved policies and
unresolved implementation mechanisms follow Sections~\ref{sec:business-rules}--\ref{sec:solution-architecture}.

\identifier{SEC-08} requires an integrity digest for each accepted document source, stored with
protected metadata and verified before association. Keep client confirmation and server acceptance timestamps distinct. A unique document reference permits one current accepted DocumentVersion; replacement creates its new identity and removes old signing/rendition rows and file references consistently. Salted password hashes, session/token
verifiers, security-event logs,
deletion/restriction recovery metadata and deployment key controls are technical
persistence concerns, not additional business actors/entities. Their implementation
and approved technical retention limits follow Section~\ref{sec:non-functional-requirements}
and Section~\ref{sec:nfr-targets}.

AccessGrant has required property scope and a unique membership/operation/property
key; there is no independent campaign/parcel scope discriminator. WorkerAvailability
uses same-workspace membership/worker FKs and a unique pair. Validate invitation
scope at creation and acceptance, including revoked/archived property context.
Security events use a 90-day server-event policy; minimal operation outcomes and
change/deletion markers use their separate 90-day policies. Store originator and
accepting owner separately for reviewed recovery. These mechanisms do not add
ordinary field-change history or override the current-only document replacement policy.

\paragraph{Provisioning fields.}
Add checked system role and nullable self-referencing administrator issuer on
\texttt{app\_user}. Existing normal users migrate to standard authority; never infer
administrator status from an Owner membership. Use a deployment-level single-client
transactional guard and pending-to-active workspace/membership checks. A protected
bootstrap procedure prepares the initial administrator; normal APIs never expose
administrator self-registration. Token verifiers and password hashes remain
technical authentication persistence. No additional domain entity is needed.
